\subsection{Suffix Array + LCP de Kasai (subcadena repetida más larga)}

\graybold{Preguntas guía:}

\begin{itemize}
\item ¿Me piden la subcadena más larga que aparece al menos dos veces en un texto (permitiendo solapamiento)?
\item ¿Necesito comparar prefijos comunes entre TODAS las subcadenas del texto, y enumerarlas explícitamente sería cuadrático o peor?
\item ¿El problema se reduce a una consulta sobre prefijos comunes de sufijos (repeticiones, subcadena común, cantidad de subcadenas distintas)?
\end{itemize}

\graybold{Utilidad:} El suffix array ordena lexicográficamente todos los sufijos de una cadena, y el arreglo LCP guarda la longitud del prefijo común entre cada par de sufijos ADYACENTES en ese orden. Juntos permiten responder en tiempo casi lineal preguntas sobre repeticiones y subcadenas comunes que de forma ingenua requerirían enumerar las $\bigO{n^2}$ subcadenas: subcadena repetida más larga, cantidad de subcadenas distintas, subcadena común más larga entre varios textos, y consultas de prefijo común entre sufijos arbitrarios.

\graybold{Complejidad:} \bigO{n \log^2 n} para construir el suffix array por duplicación (log n rondas, cada una con un ordenamiento), y \bigO{n} amortizado para el arreglo LCP con el algoritmo de Kasai.

\graybold{Fundamento teórico:} Toda subcadena de $S$ es un PREFIJO de algún sufijo de $S$, así que una subcadena que aparece dos veces corresponde a un prefijo común de dos sufijos distintos, y su longitud máxima es $\max_{i\neq j} \mathrm{lcp}(\text{suf}_i, \text{suf}_j)$. La clave es que no hace falta examinar los $\bigO{n^2}$ pares: al ordenar los sufijos lexicográficamente, la función $\mathrm{lcp}$ cumple $\mathrm{lcp}(\text{sa}_i,\text{sa}_j) = \min_{i < t \le j} \mathrm{lcp}(\text{sa}_{t-1},\text{sa}_t)$ para $i<j$ (un mínimo sobre un rango del arreglo LCP), porque un prefijo común a dos sufijos lo es también de todos los sufijos intercalados entre ellos en el orden. Como un mínimo sobre un rango nunca supera a ninguno de sus elementos, el máximo global sobre todos los pares se alcanza necesariamente en un par ADYACENTE, y la respuesta es simplemente $\max(\text{lcp})$. La construcción del suffix array por duplicación se apoya en que, conociendo el rango de cada sufijo por sus primeros $k$ caracteres, el rango por $2k$ caracteres se obtiene ordenando por el par $(\mathrm{rank}_k[i],\ \mathrm{rank}_k[i+k])$ — de ahí las \ceil{\log n} rondas. El algoritmo de Kasai calcula los LCP en orden de POSICIÓN en el texto (no de rango) usando el invariante de que, si $\mathrm{lcp}$ del sufijo que empieza en $i$ vale $h$, entonces el del sufijo que empieza en $i+1$ vale al menos $h-1$: al quitar el primer carácter de dos cadenas que comparten $h$ caracteres, siguen compartiendo $h-1$. Por eso $h$ solo decrece en una unidad por iteración pero puede crecer arbitrariamente, dando un total amortizado de \bigO{n}.

\graybold{Ejercicio aplicado}

Dados $T$ textos, para cada uno hallar la longitud de la subcadena más larga que aparece al menos dos veces en él (las dos apariciones pueden solaparse).

\graybold{I/O:} $T$ casos, cada uno una cadena en su propia línea con $|S| \le 5000$ sobre el alfabeto minúsculo $\Rightarrow$ como las cadenas no contienen espacios, se usa lectura total + tokenización (patrón 2) y cada token es directamente una cadena completa. Se atiende aparte el caso $|S| < 2$ (respuesta $0$: no puede haber repetición). Verificado contra el caso de ejemplo y contrastado con una fuerza bruta (enumerar todas las subcadenas de cada longitud en un conjunto) en 158 casos, incluyendo alfabeto binario, cadenas de un solo carácter repetido y otros casos borde, sin discrepancias. El enunciado no acota $T$, así que se midió el peor caso razonable: una cadena de $|S|=5000$ tarda \aprox{}0.02s (tanto aleatoria como en el peor caso para Kasai, \texttt{"a"} repetida 5000 veces), y $T=100$ cadenas de longitud máxima \aprox{}1.04s.

\graybold{Código solución}

\begin{lstlisting}[language=Python]
import sys

def suffix_array(s):
    n = len(s)
    # ronda inicial: ordenar los sufijos por su primer caracter
    sa = sorted(range(n), key=lambda i: s[i])
    rank = [0]*n
    r = 0
    prev = s[sa[0]]
    for idx, i in enumerate(sa):
        if s[i] != prev:
            r += 1
            prev = s[i]
        rank[i] = r
    k = 1
    # duplicacion: si ya se conoce el rango por k caracteres,
    # ordenar por el par (rank[i], rank[i+k]) da el rango por 2k
    while r < n - 1:                 # corta apenas todos los rangos son distintos
        key = [0]*n
        for i in range(n):
            # el par se empaqueta en un solo entero para ordenar mas rapido
            key[i] = rank[i] * (n + 1) + (rank[i+k] + 1 if i + k < n else 0)
        sa.sort(key=lambda i: key[i])
        new_rank = [0]*n
        r = 0
        prev = key[sa[0]]
        for i in sa:
            if key[i] != prev:
                r += 1
                prev = key[i]
            new_rank[i] = r
        rank = new_rank
        k <<= 1
    return sa, rank

def kasai(s, sa, rank):
    n = len(s)
    lcp = [0]*n
    h = 0
    # se recorre por POSICION en el texto: al avanzar de i a i+1,
    # el lcp baja como mucho en 1, asi que h nunca se reinicia a 0
    for i in range(n):
        if rank[i] > 0:
            j = sa[rank[i]-1]        # sufijo anterior en orden lexicografico
            while i + h < n and j + h < n and s[i+h] == s[j+h]:
                h += 1
            lcp[rank[i]] = h
            if h:
                h -= 1
        else:
            h = 0
    return lcp

def solve():
    data = sys.stdin.buffer.read().split()
    t = int(data[0])
    out = []
    for i in range(1, t+1):
        s = data[i].decode()
        n = len(s)
        if n < 2:
            out.append("0")
            continue
        sa, rank = suffix_array(s)
        lcp = kasai(s, sa, rank)
        # el maximo lcp sobre pares adyacentes ya es el maximo global
        out.append(str(max(lcp)))
    sys.stdout.write("\n".join(out) + "\n")

solve()
\end{lstlisting}

\graybold{Extras:} variantes de \texttt{solve()} para cuando el problema pide, además de la longitud, QUÉ subcadena es y DÓNDE aparece (\texttt{lcp\_detailed}: todas las subcadenas de longitud máxima con sus posiciones) o cuál es la subcadena más larga que aparece el MAYOR número de veces (\texttt{lcp\_repeated}, con mínimo en ventana deslizante sobre el arreglo LCP). Se llaman con \texttt{(s, sa, lcp, n)} en lugar de \texttt{max(lcp)}.

\begin{lstlisting}[language=Python]
from collections import deque, Counter

def blocks(s, sa, lcp, n, L):
    # bloques maximales con lcp >= L -> (subcadena, ocurrencias)
    res = []
    p = 1
    while p < n:
        if lcp[p] >= L:
            ini = p
            while p < n and lcp[p] >= L:
                p += 1
            occ = sorted(sa[ini-1:p])
            res.append((s[sa[ini]:sa[ini]+L], occ))
        else:
            p += 1
    return res

def lcp_detailed(s, sa, lcp, n):
    L = max(lcp)
    if L == 0:
        return "0"
    partes = [f"'{sub}' x{len(occ)} en {occ}" for sub, occ in blocks(s, sa, lcp, n, L)]
    return f"{L} -> " + " | ".join(partes)

def lcp_repeated(s, sa, lcp, n):
    # ninguna subcadena se repite mas que su primer caracter
    c = max(Counter(s).values())
    if c < 2:
        return "0"
    w = c - 1                      # ventana de w valores de lcp = c sufijos
    dq = deque()                   # minimo de ventana deslizante
    mejorL = mejorA = -1
    for p in range(1, n):
        while dq and lcp[dq[-1]] >= lcp[p]:
            dq.pop()
        dq.append(p)
        if dq[0] < p - w + 1:
            dq.popleft()
        if p >= w and lcp[dq[0]] > mejorL:
            mejorL = lcp[dq[0]]
            mejorA = p - w + 1
    occ = sorted(sa[mejorA-1:mejorA+w])
    return f"{mejorL} -> '{s[sa[mejorA]:sa[mejorA]+mejorL]}' x{len(occ)} en {occ}"
\end{lstlisting}
